\documentclass[conference]{IEEEtran}
\usepackage{cite}
\usepackage{amsmath,amssymb}
\usepackage{graphicx}
\usepackage{booktabs}
\usepackage{url}
\urlstyle{same}
\usepackage[hidelinks,breaklinks]{hyperref}
\usepackage{balance}

\begin{document}

\title{MeasureMatch-Bench: LLM Entity Resolution and Retrieval over Unstructured Hardware Measurement Corpora}

\author{\IEEEauthorblockN{Piyush Jagadish Bag}
\IEEEauthorblockA{\textit{Independent Researcher}\\
\textit{ECAD validation and test automation}\\
Santa Clara, CA, USA\\
piyushbag4@gmail.com}}

\maketitle

\begin{abstract}
Post-silicon validation teams archive instrument exports under inconsistent product aliases, heterogeneous file formats, and scenario-specific naming conventions.
Retrieving a single metric (e.g., wall power at line rate for a given SKU) can require days of manual search across thousands of files.
Existing power benchmarks target chip or LLM-inference energy, not lab measurement corpora; entity-matching research rarely addresses hardware validation aliases at archive scale.
We introduce \textbf{MeasureMatch-Bench}, a public benchmark of 24 synthetic corpora ($\sim$6{,}000 files, 46 scenario types spanning PSU, thermal, fan, throughput, and environmental measurements) with gold labels for four tasks: (T1) product identity resolution, (T2) format detection and normalization, (T3) scenario-targeted retrieval, and (T4) cross-corpus question answering.
We release harness v0.1.0 with a rule-and-fuzzy baseline (R0) that achieves 100\% T1 accuracy and 87.2\% T3 retrieval@1 on the test split, establishing a reproducible floor for embedding and open-weight LLM pipelines (R1--R2).
\end{abstract}

\begin{IEEEkeywords}
hardware validation, entity resolution, measurement data, retrieval-augmented generation, benchmark
\end{IEEEkeywords}

\section{Introduction}
Hardware validation organizations accumulate instrument exports across hundreds of product identifiers and dozens of measurement scenarios (PSU rails, thermal deltas, fan curves, line-rate throughput).
Product IDs appear under inconsistent aliases in filenames, CSV headers, and plain-text key--value blocks.
Engineers routinely spend \emph{days} manually searching shared drives and email archives to answer structured queries such as ``what was wall power at line rate for SKU~X?''

EDA flows and lab instruments capture measurements, but the \emph{corpus integration} layer---aligning aliases, normalizing formats, and retrieving scenario-specific slices---remains ad hoc.
Recent LLM work on hardware targets RTL generation and assertion mining~\cite{cvdp2025,vert2025,fixme2026}, while power benchmarks such as TokenPowerBench~\cite{tokenpowerbench2025} measure \emph{LLM inference} energy, not post-silicon lab archives.
Aerospace case studies combine LLMs with knowledge graphs for test-document extraction~\cite{desantis2024llmkg}, yet no open benchmark evaluates entity resolution and retrieval over validation-scale measurement corpora.

\textbf{Contributions.}
\begin{itemize}
  \item \textbf{MeasureMatch-Bench} --- 24 corpora, $\sim$6{,}000 synthetic measurement files, 46 scenario types, gold labels for T1--T4.
  \item \textbf{Reproducible harness} (v0.1.0) with rule baseline R0 and extension points for embedding (R1) and Qwen2.5-3B-Instruct (R2) pipelines.
  \item \textbf{Empirical floor} --- R0 scores on public splits (Table~\ref{tab:r0}) from harness output on release corpus v1.0.0.
\end{itemize}

\section{Related Work}
Table~\ref{tab:diff} summarizes differentiation. TokenPowerBench and GPU power studies address accelerator telemetry~\cite{tokenpowerbench2025,caraml2024}.
Hardware LLM benchmarks (CVDP, FIXME, VERT) target HDL and verification artifacts~\cite{cvdp2025,fixme2026,vert2025}.
General entity matching (AnyMatch) optimizes product-table linkage~\cite{anymatch2024}, not instrument-export alias chaos.
IEA-Plugin~\cite{ieaplugin2025} agents front-end industrial test analytics without a public corpus.

\begin{table}[t]
\centering
\caption{Differentiation vs.\ adjacent work}
\label{tab:diff}
\resizebox{\columnwidth}{!}{%
\begin{tabular}{@{}llll@{}}
\toprule
Prior work & Task & Data modality & Gap \\
\midrule
TokenPowerBench & LLM inference power & GPU telemetry & Not lab measurement files \\
CVDP / FIXME / VERT & RTL / verification LLM & HDL, assertions & Not measurement corpora \\
De Santis et al.~\cite{desantis2024llmkg} & LLM+KG test docs & Aerospace boards & Case study, not open bench \\
AnyMatch~\cite{anymatch2024} & Zero-shot entity match & Product tables & Not HW validation aliases \\
IEA-Plugin~\cite{ieaplugin2025} & Agent test analytics & Industrial backend & No public harness \\
\textbf{MeasureMatch-Bench} & Identity + retrieval + QA & Measurement exports & \textbf{Open multi-task bench} \\
\bottomrule
\end{tabular}}
\end{table}

\section{Benchmark Design}
\subsection{Corpus}
Each \emph{corpus} simulates one validation program archive: 40--80 canonical SKUs, 2--5 aliases each, and 250 heterogeneous CSV/TXT exports.
Twenty-four corpora (dev: 4, test: 20) total $\sim$6{,}000 files.
\textbf{46 scenario types} cover PSU (8), thermal (10), fan (10), throughput (12), and environmental (6) measurements.

\subsection{Tasks}
\begin{itemize}
  \item \textbf{T1 Identity} --- map alias string to canonical product ID.
  \item \textbf{T2 Normalize} --- detect format and emit unified schema (\texttt{format}, \texttt{scenario}, \texttt{canonical\_id}, \texttt{value}).
  \item \textbf{T3 Retrieve} --- given (\texttt{product}, \texttt{scenario}), return file path and numeric value.
  \item \textbf{T4 QA} --- answer templated questions (derived from T3 gold) with cited values.
\end{itemize}

\subsection{Metrics}
T1: accuracy; T2: record match rate; T3: retrieval@1 and value tolerance ($\pm 0.05$ absolute or 0.1\% relative); T4: value score.

\section{Reference Pipelines}
\textbf{R0 (rules + fuzzy match).}
YAML-driven alias normalization (prefix strip, PN$\rightarrow$SKU, case fold, \texttt{difflib} fuzzy cutoff 0.82) plus content/filename parsers for CSV and TXT exports.
\textbf{R1 (planned)} --- sentence embeddings + dense retrieval.
\textbf{R2 (planned)} --- Qwen2.5-3B-Instruct zero-shot prompts per task; local-only for reproducibility.

\section{Experiments}
All numbers from \texttt{eval/run\_eval.py} v0.1.0 on corpus v1.0.0 (GitHub \url{https://github.com/piyushbag/measure-match-bench}).

\begin{table}[t]
\centering
\caption{R0 aggregate scores (harness v0.1.0, 2026-07-04)}
\label{tab:r0}
\begin{tabular}{@{}lrrrrr@{}}
\toprule
Split & Corpora & T1 acc & T2 match & T3 R@1 & T4 val \\
\midrule
dev  & 4  & 1.000 & 0.823 & 0.861 & 0.952 \\
test & 20 & 1.000 & 0.819 & 0.872 & 0.959 \\
\bottomrule
\end{tabular}
\end{table}

R0 resolves all alias queries (T1) on synthetic corpora where mutations follow documented rules; T3 misses ($\sim$13\%) stem from ambiguous filename collisions when multiple aliases share scenario suffixes---motivation for R1/R2 semantic retrieval.
T2 gaps reflect canonical-ID resolution failures on noisy alias fields inside file bodies.

\section{Discussion}
MeasureMatch-Bench is intentionally synthetic: it enables public release without proprietary exports while preserving alias noise patterns seen in validation archives.
Limitations include fixed scenario taxonomy, two file formats, and R0-centric initial results; R1/R2 comparisons are future work on the same harness.
We target ITC~2027 Track~B (AI for test) with DATE~2027 as backup after arXiv archival.

\section{Conclusion}
We presented MeasureMatch-Bench, the first open benchmark for product identity alignment and scenario retrieval over hardware validation measurement corpora, with a reproducible harness and rule baseline on 24 corpora.
Future releases will add Zenodo DOI pinning, R1/R2 pipelines, and harder OOV alias splits.

\section*{Data Availability}
Corpus, gold labels, and harness: \url{https://github.com/piyushbag/measure-match-bench} (tag \texttt{v1.0.0} at release).
Zenodo DOI will be minted at v1.0.0 publication.

\bibliographystyle{IEEEtran}
\bibliography{references}
\balance
\end{document}
